\documentclass[11pt]{article}
\usepackage[margin=1in]{geometry}
\usepackage{amsmath,amssymb,amsthm,mathtools}
\usepackage{enumitem}
\usepackage{microtype}
\input{glyphtounicode}
\pdfgentounicode=1
\pdfglyphtounicode{mu}{03BC}
\pdfglyphtounicode{bardbl}{2016}
\pdfglyphtounicode{parenleftbig}{0028}
\pdfglyphtounicode{parenrightbig}{0029}
\pdfglyphtounicode{parenleftbigg}{0028}
\pdfglyphtounicode{parenrightbigg}{0029}
\pdfglyphtounicode{parenlefttp}{239B}
\pdfglyphtounicode{parenleftbt}{239D}
\pdfglyphtounicode{parenrighttp}{239E}
\pdfglyphtounicode{parenrightbt}{23A0}
\pdfglyphtounicode{circleplusdisplay}{2A01}
\pdfglyphtounicode{summationdisplay}{2211}
\pdfglyphtounicode{productdisplay}{220F}
\usepackage[colorlinks=true,linkcolor=blue,citecolor=blue,urlcolor=blue]{hyperref}
\usepackage{accsupp}
\NewCommandCopy{\OriginalHookrightarrow}{\hookrightarrow}
\NewCommandCopy{\OriginalLongmapsto}{\longmapsto}
\NewCommandCopy{\OriginalUpbracefill}{\upbracefill}
\renewcommand{\hookrightarrow}{\mathrel{
  \BeginAccSupp{method=hex,unicode,ActualText=21AA}
  \OriginalHookrightarrow\EndAccSupp{}}}
\renewcommand{\longmapsto}{\mathrel{
  \BeginAccSupp{method=hex,unicode,ActualText=27FC}
  \OriginalLongmapsto\EndAccSupp{}}}
\renewcommand{\upbracefill}{
  \BeginAccSupp{method=hex,unicode,ActualText=23DF}
  \OriginalUpbracefill\EndAccSupp{}}
\hypersetup{pdftitle={Quadratic stabilization of the canonical Foulkes--Howe map},pdfauthor={Rewritten for readability from the OpenAI manuscript of the same title}}
\newtheorem{theorem}{Theorem}
\newtheorem{corollary}[theorem]{Corollary}
\newtheorem{proposition}[theorem]{Proposition}
\newtheorem{lemma}[theorem]{Lemma}
\theoremstyle{remark}
\newtheorem{remark}[theorem]{Remark}
\DeclareMathOperator{\Sym}{Sym}
\DeclareMathOperator{\Av}{Av}
\newcommand{\C}{\mathbb C}
\newcommand{\Sa}{\mathfrak S_a}
\newcommand{\D}[2]{D_{#1,#2}}
\newcommand{\unverified}[1]{\par\smallskip\noindent\fbox{\parbox{0.95\linewidth}{\textbf{Unverified step.} #1}}\par\smallskip}
\title{Quadratic stabilization of the canonical Foulkes--Howe map}
\author{Rewritten for readability from the OpenAI manuscript of the same title}
\date{September 25, 2026}
\begin{document}
\maketitle
\begin{abstract}
For every finite-dimensional complex vector space $V$, we prove that the
canonical Foulkes--Howe map
$\Sym^b(\Sym^a V)\longrightarrow\Sym^a(\Sym^b V)$ is surjective whenever
$a\ge2$ and $b\ge a(a-1)$. This gives a quadratic stabilization bound
independent of $\dim V$.
\end{abstract}

\section{The main result}\label{sec:introduction}

Throughout, $V$ is a finite-dimensional complex vector space and
$\Sym V=\bigoplus_{k\ge0}\Sym^kV$ is its symmetric algebra. A product of
vectors always means a product in $\Sym V$.

\paragraph{Averaging.}
Let $W$ be a complex vector space and $a\ge1$ an integer. The symmetric
group $\Sa$ acts on $W^{\otimes a}$ by permuting the tensor factors. Define
the linear map $\Av\colon W^{\otimes a}\to W^{\otimes a}$ by
\begin{equation}\label{eq:average}
 \Av(w_1\otimes\cdots\otimes w_a)
 =\frac1{a!}\sum_{\pi\in\Sa}
 w_{\pi(1)}\otimes\cdots\otimes w_{\pi(a)}.
\end{equation}
We use three standard facts.
\begin{enumerate}[label=(A\arabic*),leftmargin=3.2em]
\item\label{A1} $\Av$ is a projection of $W^{\otimes a}$ onto the invariant
 tensors $(W^{\otimes a})^{\Sa}$. Its values are invariant, because
 permuting the factors only reindexes the sum. It fixes every invariant
 tensor $\tau$, because then all $a!$ summands equal $\tau$.
\item\label{A2} Let $\mathrm{pr}\colon W^{\otimes a}\to\Sym^aW$ be the
 quotient map. Since $\Av$ takes the same value on a tensor and on any
 permutation of it, $\Av$ factors as $\overline{\Av}\circ\mathrm{pr}$ with
 \[
 \overline{\Av}\colon\Sym^aW\to(W^{\otimes a})^{\Sa},\qquad
 w_1\cdots w_a\longmapsto\Av(w_1\otimes\cdots\otimes w_a).
 \]
 This map is an isomorphism with inverse $\mathrm{pr}|_{(W^{\otimes a})^{\Sa}}$.
 Indeed, $\mathrm{pr}\circ\Av=\mathrm{pr}$ gives
 $\mathrm{pr}\circ\overline{\Av}=\mathrm{id}$, and for invariant $\tau$ we
 have $\overline{\Av}(\mathrm{pr}(\tau))=\Av(\tau)=\tau$ by \ref{A1}.
\item\label{A3} $\Av(w^{\otimes a})=w^{\otimes a}$ by \ref{A1}. Hence
 $\overline{\Av}(w^a)=w^{\otimes a}$.
\end{enumerate}

\paragraph{The algebra $R$ and the map.}
Fix an integer $a\ge1$. Let $\mathcal P=(\Sym V)^{\otimes a}$ with
componentwise multiplication. (The source calls this algebra $P_a$; we
rename it because $P$ will denote a product of vectors.) The group $\Sa$
permutes the tensor factors of $\mathcal P$ by algebra automorphisms. For
each integer $j\ge0$ put
\[
 R_j=\bigl((\Sym^j V)^{\otimes a}\bigr)^{\Sa},\qquad
 R=\bigoplus_{j\ge0}R_j\subseteq\mathcal P .
\]
A product of invariant elements is invariant, $R_iR_j\subseteq R_{i+j}$,
and $\mathcal P$ is commutative. So $R$ is a graded commutative algebra.
By \ref{A2} with $W=V$ and with $W=\Sym^bV$,
\[
 R_1\cong\Sym^aV,\qquad R_b\cong\Sym^a(\Sym^bV).
\]
For every integer $b\ge0$, multiplication of $b$ elements of $R_1$ is
symmetric and multilinear, so it induces a linear map $\Sym^bR_1\to R_b$.
Under the two identifications above this is the \emph{canonical
Foulkes--Howe map}
\begin{equation}\label{eq:map}
 \mu_{a,b,V}:\Sym^b(\Sym^a V)\longrightarrow\Sym^a(\Sym^b V).
\end{equation}
Let $A\subseteq R$ be the subalgebra generated by $R_1$ and put
$A_b=A\cap R_b$. This is the span of all products of $b$ elements of $R_1$,
so $\operatorname{im}\mu_{a,b,V}=A_b$. In positive dimension,
Raicu--Sam--Weyman \cite[Section~2]{RSW} formulate this map as
$\Sym^b(D^aV)\longrightarrow D^a(\Sym^bV)$, with divided powers realized as
invariant tensors. Transporting both sides by the averaging above gives
exactly \eqref{eq:map}.

\begin{theorem}\label{thm:main}
For every integer $a\ge2$, every integer $b\ge a(a-1)$, and every
finite-dimensional complex vector space $V$, the map
$\mu_{a,b,V}$ in \eqref{eq:map} is surjective. Equivalently, $A_b=R_b$.
\end{theorem}

The case $a=1$ is the identity map in every degree. The bound $a(a-1)$ is a
sufficient bound, valid in every finite dimension at once. It is not
asserted to be sharp, and the argument does not determine the least degree
from which surjectivity holds. History and related work are collected in
Section~\ref{sec:context}; nothing there is used in the proofs.

\section{Proof at a glance}\label{sec:glance}

\paragraph{Key idea.}
Surjectivity is dualized. A linear functional on $R_b$ is the same thing as
a symmetric $a$-linear form $T$ on $\Sym^bV$. The image $A_b$ is spanned by
the tensors $P^{\otimes a}$, where $P=u_1\cdots u_b$ is a \emph{product of
$b$ vectors}. So $\mu_{a,b,V}$ is surjective if and only if
\[
 T(P,\ldots,P)=0\ \text{for all products $P$ of $b$ vectors}
 \quad\Longrightarrow\quad T=0
\]
(Proposition~\ref{prop:criterion}). This implication is proved by induction
on $a$ (Proposition~\ref{prop:vanishing}). The engine is
Lemma~\ref{lem:shift}: for two blocks of variables $U,Y$, the directional
derivative $\D YU=\sum_kY_k\,\partial/\partial U_k$ is injective on
polynomials whose $U$-degree exceeds their $Y$-degree. It is used
backwards. To show that a polynomial $h$ vanishes, apply a chain of such
derivatives until an expression already known to be zero is reached. Then
remove the derivatives one at a time by injectivity.

\paragraph{The induction step.}
Put $r=a-1$ and $m=b-r$, so $m\ge r^2$.
\begin{enumerate}[label=\textbf{Step \arabic*.},leftmargin=4.2em]
\item Test the hypothesis on $P=(x^r+ct^r)Q$, where $Q$ is a product of $m$
 vectors. This $P$ is a product of $b$ vectors because a binary form
 splits over $\C$. The coefficient of $c$ gives
 $T(x^rQ,\ldots,x^rQ,t^rQ)=0$.
\item Free the last slot: $T(x^rQ,\ldots,x^rQ,t^b)=0$. A chain of $m$
 derivatives $\D{y_i}{t}$ turns $t^b$ into a multiple of $t^rQ$, which is
 Step~1. Each is injective because the $t$-degree $b-i+1$ exceeds the
 $y_i$-degree $r$.
\item For fixed $x,t$, the form
 $(Q_1,\ldots,Q_r)\mapsto T(x^rQ_1,\ldots,x^rQ_r,t^b)$ is symmetric
 $r$-linear on $\Sym^mV$ with vanishing diagonal on products. As
 $m\ge r(r-1)$, the induction hypothesis makes it zero.
\item Remove the common factor $x^r$. A chain of $r^2$ derivatives
 $\D{x}{v_i}$ turns $T(v_1^b,\ldots,v_r^b,t^b)$ into a multiple of
 $T(x^rv_1^m,\ldots,x^rv_r^m,t^b)$, which is zero by Step~3. All $r^2$
 derivatives raise the degree of the \emph{same} block $x$, up to
 $r^2-1$. The source degree stays at least $m+1\ge r^2+1$. This
 accumulation is where the quadratic bound is needed. So $T$ vanishes on
 pure powers, and pure powers span $\Sym^bV$.
\end{enumerate}

\paragraph{Dependencies.}
\begin{itemize}[leftmargin=2em]
\item Lemma~\ref{lem:powers} (pure powers span $\Sym^kV$) feeds
 Proposition~\ref{prop:criterion}, and also the base case and the last
 line of Proposition~\ref{prop:vanishing}.
\item Lemma~\ref{lem:shift} (differential injectivity), in the chain form
 of Remark~\ref{rem:chain}, feeds Steps~2 and~4 of
 Proposition~\ref{prop:vanishing}.
\item Proposition~\ref{prop:criterion} and Proposition~\ref{prop:vanishing}
 together give Theorem~\ref{thm:main}. Theorem~\ref{thm:main} gives
 Corollary~\ref{cor:comparison}.
\end{itemize}
Source numbering: Theorem~\ref{thm:main}, Lemma~\ref{lem:shift},
Proposition~\ref{prop:vanishing} and Corollary~\ref{cor:comparison} are the
source's Theorem~1, Lemma~2, Proposition~3 and Corollary~4.
Proposition~\ref{prop:criterion} is the source's unnumbered criterion.
Lemma~\ref{lem:powers} and Remark~\ref{rem:chain} are stated in the
source's running text.

\section{Conventions and two elementary facts}\label{sec:conventions}

\paragraph{Products.}
For $k\ge1$, a \emph{product of $k$ vectors} is an element
$u_1\cdots u_k\in\Sym^kV$ with $u_i\in V$. The vectors may repeat or be
zero. A \emph{pure power} is an element $u^k$.

\begin{lemma}[pure powers span]\label{lem:powers}
For every integer $k\ge1$, the pure powers $u^k$, $u\in V$, span $\Sym^kV$.
\end{lemma}
This is used with $k=a$ in Proposition~\ref{prop:criterion} and with $k=b$
in Proposition~\ref{prop:vanishing}.
\begin{proof}
We prove the polarization identity
\begin{equation}\label{eq:polar}
 k!\,u_1\cdots u_k
 =\sum_{J\subseteq\{1,\ldots,k\}}(-1)^{k-|J|}
 \Bigl(\sum_{j\in J}u_j\Bigr)^k
 \qquad(u_1,\ldots,u_k\in V).
\end{equation}
First take the $u_i$ to be indeterminates $z_1,\ldots,z_k$ of a polynomial
ring. Let $\gamma\in\mathbb Z_{\ge0}^k$ with $|\gamma|=k$, and let $S$ be
the set of indices $i$ with $\gamma_i>0$. The monomial $z^\gamma$ occurs in
$(\sum_{j\in J}z_j)^k$ exactly when $J\supseteq S$, and then with the
multinomial coefficient $\binom{k}{\gamma}$. Its coefficient on the right
of \eqref{eq:polar} is therefore
\[
 \binom k\gamma\sum_{J\supseteq S}(-1)^{k-|J|}
 =\binom k\gamma\sum_{i=0}^{k-|S|}\binom{k-|S|}{i}(-1)^{k-|S|-i}
 =\binom k\gamma(1-1)^{k-|S|}.
\]
This is $0$ unless $|S|=k$. The only $\gamma$ with $|\gamma|=k$ and
$|S|=k$ is $\gamma=(1,\ldots,1)$, where the coefficient is $k!$. So
\eqref{eq:polar} holds for indeterminates. Applying the algebra
homomorphism $\C[z_1,\ldots,z_k]\to\Sym V$, $z_i\mapsto u_i$, gives it for
vectors. Products of $k$ vectors span $\Sym^kV$ and $k!\ne0$ in $\C$, so
\eqref{eq:polar} shows that pure powers span $\Sym^kV$.
\end{proof}

\paragraph{Vector variables.}
In Sections~\ref{sec:differential} and~\ref{sec:vanishing} we assume
$n=\dim V>0$ and fix a basis $e_1,\ldots,e_n$ of $V$. A \emph{vector
variable} $U$ is a block $(U_1,\ldots,U_n)$ of independent scalar
indeterminates. We also write $U$ for the generic vector $\sum_kU_ke_k$.
Expressions such as $U^b$ or $U^rY$ are then elements of
$\Sym V\otimes\C[\text{variables}]$, a commutative algebra over the
polynomial ring. A multilinear form $T$ on $\Sym^bV$ extends to a
multilinear form over the polynomial ring, so $T$ of such expressions is a
scalar polynomial. Substituting actual vectors for the vector variables is
a ring homomorphism. It commutes with products in $\Sym V$ and with $T$.
Following the source, we use the same letter for a vector variable and for
the vector substituted into it.

For two different vector variables $U,Y$ define
\[
 \D{Y}{U}=\sum_{k=1}^n Y_k\frac{\partial}{\partial U_k},
\]
acting on scalar polynomials, and coefficientwise on polynomials with
values in $\Sym V$. It is a derivation and it kills everything not
involving $U$. It lowers the $U$-degree by one and raises the $Y$-degree by
one. We use three facts.
\begin{enumerate}[label=(D\arabic*),leftmargin=3.2em]
\item\label{D1} If $j\ge1$ and $G$ is a $\Sym V$-valued polynomial not
 involving $U$, then $\D YU(U^jG)=j\,U^{j-1}Y\,G$. Indeed,
 $\D YU(U)=\sum_kY_ke_k=Y$, so the derivation rule gives
 $\D YU(U^j)=jU^{j-1}Y$ and $\D YU(G)=0$.
\item\label{D2} Let $T$ be an $a$-linear form on $\Sym^bV$ and
 $G_1,\ldots,G_a$ be $\Sym^bV$-valued polynomials, of which only $G_s$
 involves $U$. Then
 $\D YU\,T(G_1,\ldots,G_a)=T(G_1,\ldots,\D YUG_s,\ldots,G_a)$.
 To see this, expand $G_s=\sum_\ell g_\ell\epsilon_\ell$ in a basis
 $(\epsilon_\ell)$ of $\Sym^bV$ with scalar polynomials $g_\ell$. Then
 $T(G_1,\ldots,G_a)=\sum_\ell g_\ell\,T(\ldots,\epsilon_\ell,\ldots)$, and
 the factors $T(\ldots,\epsilon_\ell,\ldots)$ do not involve $U$.
\item\label{D3} A polynomial in finitely many variables over $\C$ that
 vanishes at every point is the zero polynomial. This follows by induction
 on the number of variables from the univariate root bound.
\end{enumerate}

\section{The canonical image and diagonal vanishing}\label{sec:map}

This section turns surjectivity of $\mu_{a,b,V}$ into a vanishing problem.
The result is used once, in the proof of Theorem~\ref{thm:main}.

\begin{proposition}[diagonal criterion]\label{prop:criterion}
Let $a,b\ge1$ be integers. The map $\mu_{a,b,V}$ is surjective if and only
if every symmetric $a$-linear form $T$ on $\Sym^b V$ satisfying
$T(P,\ldots,P)=0$ for all products $P$ of $b$ vectors is zero.
\end{proposition}

\begin{proof}
\emph{Plan.} We describe $A_b=\operatorname{im}\mu_{a,b,V}$ by explicit
spanning vectors, identify $R_b^*$ with symmetric $a$-linear forms, and
compare annihilators.

\emph{Step 1: $A_b=\operatorname{span}_{\C}\{P^{\otimes a}: P\text{ is a
product of }b\text{ vectors}\}$.}
For $u\in V$ put $z(u)=u^{\otimes a}\in R_1$. By \ref{A3}, the
identification $\Sym^aV\cong R_1$ sends $u^a$ to $z(u)$. By
Lemma~\ref{lem:powers} with $k=a$, the $u^a$ span $\Sym^aV$, so the $z(u)$
span $R_1$. The space $A_b$ is spanned by products of $b$ elements of
$R_1$. By multilinearity of the product it is spanned by products of $b$
spanning elements. Componentwise multiplication gives
\[
 z(u_1)\cdots z(u_b)=(u_1\cdots u_b)^{\otimes a}.
\]

\emph{Step 2: functionals on $R_b$ are symmetric forms.}
For $\lambda\in R_b^*$ define
\[
 T_\lambda(w_1,\ldots,w_a)
 =\lambda\bigl(\Av(w_1\otimes\cdots\otimes w_a)\bigr)
 \qquad(w_i\in\Sym^bV).
\]
This is a symmetric $a$-linear form on $\Sym^bV$, and
$\lambda\mapsto T_\lambda$ is a bijection onto all such forms.
\emph{Surjective:} given a symmetric $a$-linear $T$, let $\widetilde T$ be
the linear functional on $(\Sym^bV)^{\otimes a}$ with
$\widetilde T(w_1\otimes\cdots\otimes w_a)=T(w_1,\ldots,w_a)$, and put
$\lambda=\widetilde T|_{R_b}$. Then
\[
 T_\lambda(w_1,\ldots,w_a)
 =\frac1{a!}\sum_{\pi\in\Sa}T(w_{\pi(1)},\ldots,w_{\pi(a)})
 =T(w_1,\ldots,w_a)
\]
by symmetry of $T$.
\emph{Injective:} by \ref{A1} the averages
$\Av(w_1\otimes\cdots\otimes w_a)$ span $R_b$, so $T_\lambda$ determines
$\lambda$.

\emph{Step 3: the diagonal.}
For a product $P$ of $b$ vectors, \ref{A3} gives
$T_\lambda(P,\ldots,P)=\lambda(P^{\otimes a})$.

\emph{Step 4: conclusion.}
By Steps 1 and 3, $\lambda$ vanishes on $A_b$ if and only if
$T_\lambda(P,\ldots,P)=0$ for every product $P$ of $b$ vectors. Since
$R_b$ is finite-dimensional, $A_b=R_b$ if and only if the only
$\lambda\in R_b^*$ vanishing on $A_b$ is $\lambda=0$. By Step 2 this is the
stated condition on forms $T$.
\end{proof}

\begin{remark}
This explains the reversed roles of the two parameters in the equation
interpretation. The multiplication algebra concerns products of $a$ linear
forms in coordinate degree $b$. Here $T(P,\ldots,P)$ is a homogeneous
degree-$a$ polynomial on $\Sym^b V$, tested on products of $b$ vectors,
viewed as linear forms on $V^*$. The pairing just proved connects these two
descriptions; compare the product parametrization in
\cite[Section~2]{RSW}.
\end{remark}

\section{Differential injectivity}\label{sec:differential}

The induction in Section~\ref{sec:vanishing} recovers a polynomial from a
sequence of its directional derivatives. The following lemma permits each
derivative to be cancelled. It is used in Steps~2 and~4 of
Proposition~\ref{prop:vanishing}. Recall $n=\dim V>0$ and the operator
$\D YU$ from Section~\ref{sec:conventions}.

The criterion is the positive-weight injectivity in the usual
$\mathfrak{sl}_2$ weight strings; compare
\cite[Section~11.1, Claim~11.4 and equation~(11.5)]{FultonHarris}.
Ikenmeyer's argument \cite[Section~5, Claim~5.1]{Ikenmeyer} decomposes
tensors according to untouched positions before applying a two-coordinate
weight shift. The proof below is a direct inner-product proof with all
additional variables present.

\begin{lemma}\label{lem:shift}
Let $p,q$ be nonnegative integers. On polynomials homogeneous of degree $p$
in $U$ and degree $q$ in $Y$, the operator $\D{Y}{U}$ is injective if
$p>q$. The polynomials may depend on any finite collection of additional
variables independent of $U$ and $Y$.
\end{lemma}

\begin{proof}
\emph{Plan.} Put an inner product on the polynomial ring for which the
adjoint of $D=\D YU$ is the reversed operator $\D UY$. The commutator of
the two operators acts on the given bidegree as the scalar $p-q$. This
gives a norm identity that forces $h=0$ when $Dh=0$.

\emph{Notation.} Let $X=(X_1,\ldots,X_N)$ list all scalar variables: the
coordinates of $U$, of $Y$, and the additional variables. Write
$D=\D{Y}{U}$ and $D'=\D{U}{Y}$. (The source calls the second operator $E$.)

\emph{Step 1: the inner product.}
The source says: give the polynomial ring the positive definite Fischer, or
apolar, Hermitian inner product
\cite[Section~2, equation~(8) and Proposition~3]{AldazRender}, in which the
monomials are orthonormal, $\|X^\alpha\|^2=1$.

\unverified{The normalization $\|X^\alpha\|^2=1$ does not support Step~2.
If the monomials are orthonormal, then
$\langle\partial_{X_j}X_j^2,\,X_j\rangle=2$ but
$\langle X_j^2,\,X_j\cdot X_j\rangle=1$, so differentiation is not adjoint
to multiplication, and identity \eqref{eq:norm} fails (see the sanity check
after the proof). The inner product usually called Fischer or apolar has
orthogonal monomials with $\|X^\alpha\|^2=\alpha!$. With that normalization
every remaining step checks, and the statement of the lemma is unaffected.}

For the rest of the proof, $\langle\cdot,\cdot\rangle$ therefore denotes
the Hermitian inner product
\[
 \Bigl\langle\sum_\alpha c_\alpha X^\alpha,\ \sum_\alpha d_\alpha X^\alpha\Bigr\rangle
 =\sum_\alpha\alpha!\,c_\alpha\overline{d_\alpha},
 \qquad \alpha!=\alpha_1!\cdots\alpha_N!.
\]
It is positive definite because
$\langle h,h\rangle=\sum_\alpha\alpha!\,|c_\alpha|^2$.

\emph{Step 2: differentiation is adjoint to multiplication.}
For monomials and any index $j$,
\[
 \langle\partial_{X_j}X^\alpha,\,X^\beta\rangle
 =\alpha_j\,(\alpha-e_j)!\,\delta_{\alpha-e_j,\beta}
 =\alpha!\,\delta_{\alpha,\beta+e_j}
 =\langle X^\alpha,\,X_jX^\beta\rangle .
\]
Here $e_j$ is the $j$-th unit multi-index, and both sides are $0$ if
$\alpha_j=0$. By sesquilinearity,
$\langle\partial_{X_j}g,h\rangle=\langle g,X_jh\rangle$ for all polynomials
$g,h$. Taking complex conjugates and exchanging the roles of $g$ and $h$
gives $\langle X_jg,h\rangle=\langle g,\partial_{X_j}h\rangle$. Applying
these two rules to $D=\sum_kY_k\,\partial_{U_k}$ gives
\begin{equation}\label{eq:adjoint}
 \langle Dg,h\rangle=\langle g,D'h\rangle,
 \qquad D'=\sum_kU_k\,\partial_{Y_k}.
\end{equation}

\emph{Step 3: the commutator.}
For indices $j,k$, using
$\partial_{Y_j}Y_k=\delta_{jk}+Y_k\partial_{Y_j}$ and
$\partial_{U_k}U_j=\delta_{jk}+U_j\partial_{U_k}$,
\begin{align*}
 (U_j\partial_{Y_j})(Y_k\partial_{U_k})
 &=\delta_{jk}\,U_j\partial_{U_k}+U_jY_k\,\partial_{Y_j}\partial_{U_k},\\
 (Y_k\partial_{U_k})(U_j\partial_{Y_j})
 &=\delta_{jk}\,Y_k\partial_{Y_j}+Y_kU_j\,\partial_{U_k}\partial_{Y_j}.
\end{align*}
Subtracting and summing over $j,k$ gives
\[
 D'D-DD'=N_U-N_Y,
 \qquad
 N_U=\sum_k U_k\partial_{U_k},\quad
 N_Y=\sum_k Y_k\partial_{Y_k}.
\]
By Euler's identity, $N_Uh=ph$ and $N_Yh=qh$ when $h$ is homogeneous of
degree $p$ in $U$ and $q$ in $Y$.

\emph{Step 4: the norm identity.}
Let $h$ have the indicated bidegree. By \eqref{eq:adjoint},
$\|Dh\|^2=\langle h,D'Dh\rangle$, which is real and so equals
$\langle D'Dh,h\rangle$. Also by \eqref{eq:adjoint},
$\|D'h\|^2=\langle DD'h,h\rangle$. Hence, by Step 3,
\begin{equation}\label{eq:norm}
 \|Dh\|^2-\|D'h\|^2=\langle(D'D-DD')h,h\rangle=(p-q)\|h\|^2.
\end{equation}

\emph{Step 5: conclusion.}
If $Dh=0$ and $p>q$, the left side of \eqref{eq:norm} is
$-\|D'h\|^2\le0$ and the right side is $\ge0$. So both are $0$, and
$(p-q)\|h\|^2=0$ forces $h=0$.
\end{proof}

\emph{Sanity check of \eqref{eq:norm}.} Take $n=1$ and $h=U_1^2$, so
$(p,q)=(2,0)$, $Dh=2U_1Y_1$ and $D'h=0$. With $\|X^\alpha\|^2=\alpha!$ we
get $\|Dh\|^2=4$ and $(p-q)\|h\|^2=2\cdot2=4$. With orthonormal monomials
we would get $4$ on the left and $2\cdot1=2$ on the right.

\begin{remark}[chains of operators]\label{rem:chain}
Lemma~\ref{lem:shift} is always applied in the following form. Let
$D_1,\ldots,D_s$ be operators of the type $\D{Y}{U}$, possibly with
different blocks, and let $h$ be a polynomial. Put
$h_i=D_i\cdots D_1h$, with $h_0=h$. Suppose that for each $i=1,\ldots,s$
the polynomial $h_{i-1}$ is homogeneous in the source block of $D_i$ of
some degree $p_i$ and in the target block of $D_i$ of some degree $q_i$,
with $p_i>q_i$. Then $h_s=0$ implies $h=0$. Indeed,
$D_sh_{s-1}=0$ gives $h_{s-1}=0$ by the lemma, then $h_{s-2}=0$, and so
on. All blocks other than the two belonging to $D_i$ count as additional
variables.
\end{remark}

\section{Vanishing on products}\label{sec:vanishing}

We now prove the vanishing statement required by
Proposition~\ref{prop:criterion}.

\begin{proposition}\label{prop:vanishing}
Let $V$ be a finite-dimensional complex vector space, and let $a\ge1$ and
$b\ge1$ be integers satisfying $b\ge a(a-1)$.
If a symmetric $a$-linear form $T$ on $\Sym^b V$ satisfies
\begin{equation}\label{eq:diagonal}
 T(P,\ldots,P)=0
 \qquad\text{for every product $P$ of $b$ vectors},
\end{equation}
then $T=0$.
\end{proposition}

\begin{proof}
\emph{Plan.} Induct on $a$. In the induction step, Steps~1 and~2 show that
$T(x^rQ,\ldots,x^rQ,t^b)=0$, with $r=a-1$ slots carrying a common product
and one slot free. Step~3 applies the induction hypothesis to those $r$
slots. Step~4 removes the common factor $x^r$ and concludes.

\paragraph{Trivial case and base case.}
If $V=0$ then $\Sym^bV=0$ because $b\ge1$, so $T=0$. Suppose
$n=\dim V>0$. We induct on $a$, proving the assertion for all admissible
$b$ and all $V$ at once. For $a=1$, $T$ is a linear form vanishing on all
products of $b$ vectors. These span $\Sym^bV$, so $T=0$.

\paragraph{Induction step: notation.}
Let $a\ge2$, assume the assertion for $a-1$, and put
\begin{equation}\label{eq:rm}
 r=a-1,\qquad m=b-r.
\end{equation}
The hypothesis $b\ge a(a-1)=r(r+1)$ gives
\begin{equation}\label{eq:range}
 b\ge r(r+1),\qquad m\ge r^2\ge1.
\end{equation}

\paragraph{Step 1: a split binary form.}
\emph{Claim:} for all vectors $x,t,y_1,\ldots,y_m\in V$, with
$Q=y_1\cdots y_m$,
\begin{equation}\label{eq:binary}
 T(\underbrace{x^rQ,\ldots,x^rQ}_{r\text{ slots}},t^rQ)=0.
\end{equation}

\emph{(1a) $P_c=(x^r+ct^r)Q$ is a product of $b$ vectors for every
$c\in\C$.}
Let $\xi,\eta$ be indeterminates. (The source writes $X,Y$, and writes
$\alpha$ for the number called $\beta$ here.) If $c\ne0$, choose
$\beta\in\C$ with $\beta^r=-c$ and put $\zeta=e^{2\pi i/r}$. Since
$\prod_{j=0}^{r-1}(\xi-\zeta^j\omega)=\xi^r-\omega^r$, taking
$\omega=\beta\eta$ gives
\[
 \xi^r+c\eta^r=\prod_{j=0}^{r-1}(\xi-\beta\zeta^j\eta).
\]
If $c=0$ the factorization is $\xi^r$. Apply the algebra homomorphism
$\C[\xi,\eta]\to\Sym V$, $\xi\mapsto x$, $\eta\mapsto t$. It exists for
any $x,t$; they need not be independent. This shows that $x^r+ct^r$ is a
product of $r$ vectors $x-\beta\zeta^jt$. Multiplying by $Q$ gives a
product of $r+m=b$ vectors.

\emph{(1b) Extracting the coefficient of $c$.}
Write $P_c=x^rQ+c\,t^rQ$. By multilinearity and symmetry of $T$,
\[
 T(P_c,\ldots,P_c)=\sum_{k=0}^{a}\binom ak c^k\,
 T(\underbrace{x^rQ,\ldots,x^rQ}_{a-k},\underbrace{t^rQ,\ldots,t^rQ}_{k}).
\]
By (1a) and \eqref{eq:diagonal} this polynomial in $c$ vanishes for every
$c\in\C$, so each coefficient is $0$. The coefficient of $c$ is
$a\,T(x^rQ,\ldots,x^rQ,t^rQ)$ with $a-1=r$ slots $x^rQ$, and $a\ne0$.
This is \eqref{eq:binary}.

\paragraph{Step 2: freeing the last slot.}
\emph{Claim:} for all $x,t\in V$ and every product $Q$ of $m$ vectors,
\begin{equation}\label{eq:free}
 T(\underbrace{x^rQ,\ldots,x^rQ}_{r\text{ slots}},t^b)=0.
\end{equation}
\emph{Idea.} Differentiation in $t$ inserts the factors of $Q$ into the
last slot one at a time and ends at \eqref{eq:binary}.
Remark~\ref{rem:chain} then recovers the vanishing before the insertions.

Fix a vector $x\in V$. Let $t,y_1,\ldots,y_m$ be independent vector
variables, $Q=y_1\cdots y_m$, and
\[
 H_x=T(\underbrace{x^rQ,\ldots,x^rQ}_{r\text{ slots}},t^b),
\]
a scalar polynomial in the coordinates of $t,y_1,\ldots,y_m$.

\emph{(2a) Formula.} For $0\le k\le m$,
\begin{equation}\label{eq:first-composition}
 \D{y_k}{t}\cdots\D{y_1}{t}H_x
 =\frac{b!}{(b-k)!}\,
 T(x^rQ,\ldots,x^rQ,\;t^{b-k}y_1\cdots y_k),
\end{equation}
with the empty composition for $k=0$. The proof is by induction on $k$.
Only the last slot involves $t$, so by \ref{D2} the operator $\D{y_k}{t}$
acts on the last slot alone. By \ref{D1} with $G=y_1\cdots y_{k-1}$,
\[
 \D{y_k}{t}\bigl(t^{b-k+1}y_1\cdots y_{k-1}\bigr)
 =(b-k+1)\,t^{b-k}y_1\cdots y_k .
\]
Here $b-k+1\ge b-m+1=r+1\ge1$. The scalar factors multiply to
$b(b-1)\cdots(b-k+1)=b!/(b-k)!$.

\emph{(2b) The end of the chain is zero.}
For $k=m$ we have $b-m=r$, so the right side of
\eqref{eq:first-composition} is
$\frac{b!}{r!}T(x^rQ,\ldots,x^rQ,t^rQ)$. By \eqref{eq:binary} it vanishes
for every choice of vectors $t,y_1,\ldots,y_m$. By \ref{D3} it is the zero
polynomial.

\emph{(2c) Degree check.}
Let $1\le i\le m$. By \eqref{eq:first-composition} with $k=i-1$, the
polynomial to which the $i$-th operator $\D{y_i}{t}$ is applied is
homogeneous in $t$ of degree $p=b-i+1$. In $y_i$ it is homogeneous of
degree $q=r$: the variable $y_i$ occurs once in each of the $r$ slots
$x^rQ$ and not in the last slot, since earlier operators inserted only
$y_1,\ldots,y_{i-1}$. Thus
\[
 p=b-i+1,\qquad q=r,\qquad p-q=m-i+1\ge1 .
\]

\emph{(2d) Conclusion.}
By (2b), (2c) and Remark~\ref{rem:chain}, $H_x=0$ as a polynomial.
Substituting vectors for $t,y_1,\ldots,y_m$ gives \eqref{eq:free}, since
$x$ was arbitrary.

\paragraph{Step 3: induction on the remaining slots.}
Fix $x,t\in V$ and define
\[
 S_{x,t}(Q_1,\ldots,Q_r)
 :=T(x^rQ_1,\ldots,x^rQ_r,t^b)
 \qquad(Q_i\in\Sym^mV).
\]
Multiplication by $x^r$ is a linear map $\Sym^mV\to\Sym^bV$ and $T$ is
symmetric and multilinear, so $S_{x,t}$ is a symmetric $r$-linear form on
$\Sym^mV$. By \eqref{eq:free}, $S_{x,t}(Q,\ldots,Q)=0$ for every product
$Q$ of $m$ vectors. By \eqref{eq:range}, $m\ge1$ and
$m\ge r^2\ge r(r-1)$. So the induction hypothesis applies, with $a-1=r$
in place of $a$ and $m$ in place of $b$, and gives $S_{x,t}=0$:
\begin{equation}\label{eq:inductive-vanishing}
 T(x^rQ_1,\ldots,x^rQ_r,t^b)=0
 \qquad(Q_1,\ldots,Q_r\in\Sym^m V;\ x,t\in V).
\end{equation}

\paragraph{Step 4: removing the common power.}
\emph{Claim:} $T(v_1^b,\ldots,v_r^b,t^b)=0$ for all
$v_1,\ldots,v_r,t\in V$.

\emph{Idea.} In \eqref{eq:inductive-vanishing} the $Q_i$ are independent,
and the only shared constraint left is the factor $x^r$. Differentiating
$v_i^b$ towards $x$ exactly $r$ times produces $x^rv_i^m$.

Fix a vector $t\in V$. Let $v_1,\ldots,v_r$ and $x$ be independent vector
variables; from now on $x$ is a variable, not a fixed vector. Put
\[
 f=T(v_1^b,\ldots,v_r^b,t^b).
\]
This polynomial does not involve $x$. Apply $\D{x}{v_1}$ $r$ times, then
$\D{x}{v_2}$ $r$ times, and so on up to $\D{x}{v_r}$. For $1\le i\le r$ and
$0\le j\le r$ let
\[
 f_{i,j}=\D{x}{v_i}^{\,j}\,\D{x}{v_{i-1}}^{\,r}\cdots\D{x}{v_1}^{\,r}f .
\]

\emph{(4a) Formula.}
\[
 f_{i,j}=\Bigl(\frac{b!}{m!}\Bigr)^{i-1}\frac{b!}{(b-j)!}\;
 T\bigl(x^rv_1^m,\ldots,x^rv_{i-1}^m,\;x^{j}v_i^{\,b-j},\;
 v_{i+1}^b,\ldots,v_r^b,\;t^b\bigr).
\]
The proof is by induction along the chain. Only slot $i$ involves $v_i$,
so by \ref{D2} the operator $\D{x}{v_i}$ acts on slot $i$ alone. It does
not differentiate the $x$ factors already inserted, nor any other slot. By
\ref{D1} with $G=x^{j}$,
$\D{x}{v_i}(v_i^{\,b-j}x^{j})=(b-j)\,v_i^{\,b-j-1}x^{j+1}$ for
$0\le j<r$. Also $f_{i,r}=f_{i+1,0}$ and $b-r=m$. In particular
\begin{equation}\label{eq:second-composition}
 \D{x}{v_r}^{\,r}\cdots\D{x}{v_1}^{\,r}f
 =\left(\frac{b!}{m!}\right)^r
 T(x^rv_1^m,\ldots,x^rv_r^m,t^b)=0.
\end{equation}
The last equality holds for every choice of vectors by
\eqref{eq:inductive-vanishing} with $Q_i=v_i^m$. By \ref{D3} it holds as a
polynomial identity.

\emph{(4b) Degree check.}
The operator leading from $f_{i,j-1}$ to $f_{i,j}$, for $1\le j\le r$, is
$\D{x}{v_i}$, with source block $v_i$ and target block $x$. By (4a),
$f_{i,j-1}$ is homogeneous in $v_i$ of degree $p=b-j+1$ and in $x$ of
degree $q=(i-1)r+(j-1)$. All $r^2$ operators raise the degree of the same
block $x$. Using \eqref{eq:range},
\[
 p\ \ge\ b-r+1=m+1\ \ge\ r^2+1,\qquad
 q\ \le\ (r-1)r+(r-1)=r^2-1,
\]
so $p>q$ before every operator. All other blocks are spectators.

\emph{(4c) Conclusion.}
By \eqref{eq:second-composition}, (4b) and Remark~\ref{rem:chain}, $f=0$
as a polynomial. Substituting vectors, and recalling that $t$ was
arbitrary, proves the claim: $T$ vanishes on every tuple of pure powers.

\paragraph{End of the induction step.}
Pure powers span $\Sym^bV$ by Lemma~\ref{lem:powers} with $k=b$. Since $T$
is multilinear and vanishes on all tuples of pure powers, $T=0$.
\end{proof}

\begin{proof}[Proof of Theorem~\ref{thm:main}]
Let $a\ge2$ and $b\ge a(a-1)$, so $b\ge2\ge1$.
Proposition~\ref{prop:vanishing} shows that every symmetric $a$-linear form
on $\Sym^bV$ whose diagonal vanishes on products of $b$ vectors is zero. By
Proposition~\ref{prop:criterion}, $\mu_{a,b,V}$ is surjective.
\end{proof}

\section{Representation comparison and scope}\label{sec:comparison}

\begin{corollary}\label{cor:comparison}
For every finite-dimensional complex vector space $V$, every integer
$a\ge2$, and every integer $b\ge a(a-1)$, there exists a
$\mathrm{GL}(V)$-equivariant embedding
\[
 \Sym^a(\Sym^b V)\hookrightarrow\Sym^b(\Sym^a V).
\]
In particular, this gives the sixth symmetric-power comparison for $b\ge30$.
\end{corollary}

\begin{proof}
The group $\mathrm{GL}(V)$ acts on $\Sym V$ by algebra automorphisms. So it
acts on $(\Sym V)^{\otimes a}$ by algebra automorphisms commuting with
$\Sa$. Averaging and multiplication therefore commute with the action of
$\mathrm{GL}(V)$, and the surjection $\mu_{a,b,V}$ of
Theorem~\ref{thm:main} is equivariant. Finite-dimensional polynomial
representations over $\C$ are completely reducible
\cite[Chapter~I, Appendix~A, Sections~7--8]{Macdonald}. Choose an invariant
complement $M$ to the kernel of $\mu_{a,b,V}$. The restriction
$\mu_{a,b,V}|_M$ is an equivariant isomorphism onto $\Sym^a(\Sym^b V)$. Its
inverse followed by the inclusion of $M$ in $\Sym^b(\Sym^a V)$ is the
required embedding. For $a=6$ the range is $b\ge6\cdot5=30$. This
construction asserts existence and does not identify the embedding with a
prescribed canonical map.
\end{proof}

\begin{remark}
The differential lemma holds for every finite number of coordinates and
arbitrary additional variables, while the degree comparisons depend only
on $a$ and $b$. Together with the finite-dimensional annihilator argument,
this proves the same bound in every finite dimension, without a
stable-dimension reduction. The written proof uses positivity of the
Fischer Hermitian inner product, splitting of binary forms over $\C$, and
the nonzero characteristic-zero scalar factors in the coefficient and
differentiation identities.
\end{remark}

\section{Context and related work}\label{sec:context}

This section reports the source's account of the literature. It is not
used in the proofs.

\paragraph{The stabilization question.}
For fixed $a$ and $V$, the stabilization question asks for a nonnegative
degree $b_0$ such that $\mu_{a,b,V}$ is surjective for every $b\ge b_0$.
One may ask for the least such degree in a fixed dimension, or for a bound
valid in every finite dimension at once. Theorem~\ref{thm:main} is of the
second kind.

\paragraph{History and significance.}
The affine Chow variety of decomposable $a$-forms consists of products of
$a$ linear forms in $\Sym^a(V^*)$. Its coordinate ring identifies with
$A$, whose normalization is $R$; here $b$ is the coordinate degree
\cite[Section~2, Theorem~2.3]{RSW}. This interpretation connects the
canonical map to the geometry of symmetric products; see also
Landsberg \cite[Section~7.6, Propositions~7.11--7.12]{Landsberg}.
Brion established eventual surjectivity \cite{Brion1993} and subsequently
gave an effective bound involving both $a$ and the dimension
\cite[Theorem~3.3]{Brion}. In the normalization picture, eventual
surjectivity says that the graded quotient $R/A$ has only finitely many
nonzero components. The least stabilization degree is one more than its
highest nonzero degree, or zero if $R=A$. Raicu, Sam, and Weyman bound the
regularity of the Segre ring and modules of covariants, including the
normalization \cite[Theorem~3.9]{RSW}. Their theorem does not compute the
top degree of the separate quotient $R/A$ that controls stabilization.
Theorem~\ref{thm:main} gives the sufficient bound $a(a-1)$ in every
dimension. It also answers positively the request for a polynomial bound in
\cite[Problem~7.19]{Landsberg}: in that source's notation $d=b$, $n=a$,
and $w=\dim V$, the bound is $d\ge n(n-1)$, including when $w=n$.

\paragraph{Relation to the Foulkes conjecture.}
The canonical-map question must be distinguished from the ordinary
Foulkes conjecture \cite{Foulkes1950}, which asks, for $a\le b$, for a
$\mathrm{GL}(V)$-equivariant injection
$\Sym^a(\Sym^b V)\hookrightarrow\Sym^b(\Sym^a V)$
\cite[Conjecture~2.10]{RSW}.
Surjectivity in Theorem~\ref{thm:main} implies that representation
comparison in the stated range; Corollary~\ref{cor:comparison} gives the
precise statement. Existence of an embedding alone does not specify the
canonical map. At $a=b$ the two representations are the same, so their
ordinary comparison is automatic, while a specified canonical endomorphism
can still fail to be invertible. Indeed, the canonical map has nonzero
kernels at the diagonal pairs $(a,b)=(5,5)$ and $(6,6)$ in suitable
dimensions
\cite[arXiv version, Theorem~6(b),(c) and Remark~7]{CIM}, both outside the
range of Theorem~\ref{thm:main}.

\paragraph{Relation to propagation arguments.}
Successive weight shifts also occur in Ikenmeyer's proof
\cite[Section~5]{Ikenmeyer} of McKay's propagation theorem \cite{McKay}.
In Ikenmeyer's formulation, the canonical maps in the opposite direction
$\Sym^a(\Sym^b V)\longrightarrow\Sym^b(\Sym^a V)$ are considered with
$\dim V\ge b$ for each $b$. For fixed $a$ and $B\ge a$, injectivity at
$(a,B)$ propagates to every $(a,b)$ with $b\ge B$
\cite[Theorem~1.3 and Section~3]{Ikenmeyer}.
Here the two differential sequences (Steps~2 and~4) surround an induction
on the number of slots and give a surjectivity range directly. The first
sequence uses a different target block at each step. The second raises the
degree of one shared target block at every step, and this accumulation
leads to the quadratic sufficient bound.

\section{Verification checklist}\label{sec:checklist}

\begin{enumerate}[leftmargin=2em]
\item \textbf{Inner product in Lemma~\ref{lem:shift} (Steps 1--2 of its
 proof).} The source's normalization $\|X^\alpha\|^2=1$ is flagged as an
 unverified step. Check that with $\|X^\alpha\|^2=\alpha!$ differentiation
 is adjoint to multiplication, and that nothing else in the paper depends
 on the normalization.
\item \textbf{Commutator and norm identity (Steps 3--4 of the same proof).}
 Check $D'D-DD'=N_U-N_Y$ and \eqref{eq:norm}. The sanity check after the
 proof tests the sign and the constant.
\item \textbf{Image description (Proposition~\ref{prop:criterion}, Step 1).}
 Check that $A_b$ is spanned by the tensors $P^{\otimes a}$. This uses
 Lemma~\ref{lem:powers} with $k=a$ and \ref{A3}.
\item \textbf{Step 1 of Proposition~\ref{prop:vanishing}.} Check that
 $(x^r+ct^r)Q$ is a product of exactly $r+m=b$ vectors for every $c$, and
 that the coefficient of $c$ is $a\,T(x^rQ,\ldots,x^rQ,t^rQ)$.
\item \textbf{Step 2, degree check (2c).} The $y_i$-degree before the $i$-th
 operator is $r$, not $0$, because $Q$ already sits in the first $r$
 slots. The inequality is $p-q=m-i+1\ge1$ for $i\le m$.
\item \textbf{Step 3.} The induction hypothesis is applied to $(r,m)$ in
 place of $(a,b)$. Check $m\ge1$ and $m\ge r(r-1)$ from \eqref{eq:range}.
\item \textbf{Step 4, degree check (4b).} The $x$-degree accumulates over
 all $r^2$ operators, up to $r^2-1$, while the $v_i$-degree is at least
 $m+1\ge r^2+1$. This is where $b\ge a(a-1)$ is used.
\item \textbf{Pointwise versus polynomial vanishing.} In (2b) and
 \eqref{eq:second-composition} a pointwise identity is upgraded to a
 polynomial identity by \ref{D3} before Lemma~\ref{lem:shift} is applied.
 The formulas \eqref{eq:first-composition} and (4a) rest on \ref{D1} and
 \ref{D2}.
\end{enumerate}

\begin{thebibliography}{99}
\bibitem{AldazRender}
J.~M.~Aldaz and H.~Render,
\href{https://doi.org/10.1007/s13324-023-00844-4}{\emph{A Fischer type decomposition theorem from the apolar inner product}},
Analysis and Mathematical Physics \textbf{13} (2023), article~91.

\bibitem{Brion1993}
M.~Brion,
\href{https://doi.org/10.1007/BF03026558}{\emph{Stable properties of plethysm: on two conjectures of Foulkes}},
Manuscripta Mathematica \textbf{80} (1993), no.~4, 347--371.

\bibitem{Brion}
M.~Brion,
\href{https://smf.emath.fr/publications/sur-certains-modules-gradues-associes-aux-produits-symetriques}{\emph{Sur certains modules gradu\'es associ\'es aux produits sym\'etriques}},
in \emph{Alg\`ebre non commutative, groupes quantiques et invariants}
(Reims, 1995), S\'eminaires et Congr\`es \textbf{2},
Soci\'et\'e Math\'ematique de France, Paris, 1997, 157--183.

\bibitem{CIM}
M.-W.~Cheung, C.~Ikenmeyer, and S.~Mkrtchyan,
\href{https://doi.org/10.1016/j.jsc.2016.09.002}{\emph{Symmetrizing tableaux and the 5th case of the Foulkes conjecture}},
Journal of Symbolic Computation \textbf{80} (2017), part~3, 833--843.
\href{https://arxiv.org/abs/1509.03944v1}{arXiv:1509.03944v1}.

\bibitem{Foulkes1950}
H.~O.~Foulkes,
\href{https://doi.org/10.1112/jlms/s1-25.3.205}{\emph{Concomitants of the quintic and sextic up to degree four in the coefficients of the ground form}},
Journal of the London Mathematical Society \textbf{25} (1950), no.~3,
205--209.

\bibitem{FultonHarris}
W.~Fulton and J.~Harris,
\href{https://doi.org/10.1007/978-1-4612-0979-9}{\emph{Representation Theory: A First Course}},
Graduate Texts in Mathematics \textbf{129}, Springer-Verlag, New York, 1991.

\bibitem{Ikenmeyer}
C.~Ikenmeyer,
\emph{On McKay's propagation theorem for the Foulkes conjecture},
\href{https://arxiv.org/abs/1509.04957v1}{arXiv:1509.04957v1}, 2015.

\bibitem{Landsberg}
J.~M.~Landsberg,
\href{https://doi.org/10.1007/s11565-014-0202-7}{\emph{Geometric complexity theory: an introduction for geometers}},
Annali dell'Universit\`a di Ferrara \textbf{61} (2015), no.~1, 65--117.
\href{https://arxiv.org/abs/1305.7387}{arXiv:1305.7387}.
Locators refer to the published version.

\bibitem{Macdonald}
I.~G.~Macdonald,
\href{https://doi.org/10.1093/oso/9780198534891.001.0001}{\emph{Symmetric Functions and Hall Polynomials}},
second edition, Oxford Mathematical Monographs, Clarendon Press, Oxford, 1995.

\bibitem{McKay}
T.~McKay,
\href{https://doi.org/10.1016/j.jalgebra.2007.12.003}{\emph{On plethysm conjectures of Stanley and Foulkes}},
Journal of Algebra \textbf{319} (2008), no.~5, 2050--2071.

\bibitem{RSW}
C.~Raicu, S.~V.~Sam, and J.~Weyman,
\href{https://doi.org/10.1007/s10013-021-00527-2}{\emph{On some modules supported in the Chow variety}},
Vietnam Journal of Mathematics \textbf{50} (2022), no.~2, 501--521.
\href{https://arxiv.org/abs/2108.10910v1}{arXiv:2108.10910v1}.
Locators refer to this preprint version.
\end{thebibliography}

\end{document}
